\section{Continuous functions}
\label{sec:continuous-functions}
Continuity is always relative to a function's domain. At a point
$x_0$ in that domain, Definition~17.1 requires
$f(x_n)\to f(x_0)$ for every domain sequence $x_n\to x_0$.
The equivalent epsilon--delta condition of Theorem~17.2 tests only
points $x$ in the domain as well. This qualification matters for
square roots, quotients, and isolated domain points.

Two simple starting functions are continuous on $\R$: the identity
$I(x)=x$ and every constant function $C(x)=c$. Indeed, for
$x_n\to x_0$, one has $I(x_n)=x_n\to x_0=I(x_0)$ and
$C(x_n)=c\to c=C(x_0)$. Theorems~17.3--17.5 then supply
scalar multiples, sums, products, quotients with nonzero
denominators, and compositions. Each application below checks
the domain or sign condition it needs.

\exercise{17.1}{Domains and values of two compositions}
\goal For $f(x)=\sqrt{4-x}$ on $(-\infty,4]$ and
$g(x)=x^2$ on $\R$, answer all four domain and composition questions.
\idea A sum or product needs both original values. A composition
also needs the inner value to belong to the outer function's domain.

\step{(a) Find the domains before evaluating anything.}
For $f+g$ and $fg$, intersect the two domains:
\[
 \operatorname{dom}(f+g)=\operatorname{dom}(fg)
 =(-\infty,4]\cap\R=(-\infty,4].
\]
Their formulas there are $\sqrt{4-x}+x^2$ and
$x^2\sqrt{4-x}$, respectively.

For $f\circ g$, the inner function $g$ is defined for every real
$x$, but its value must satisfy $g(x)\le4$. This means
$x^2\le4$, or equivalently $-2\le x\le2$. Hence
\[
 (f\circ g)(x)=\sqrt{4-x^2},\qquad
 \operatorname{dom}(f\circ g)=[-2,2].
\]
For $g\circ f$, first $x\le4$ is required to define $f(x)$.
The value $f(x)$ is then real, and $g$ accepts every real input.
There is no additional restriction. Squaring the nonnegative
square root gives
\[
 (g\circ f)(x)=(\sqrt{4-x})^2=4-x,\qquad
 \operatorname{dom}(g\circ f)=(-\infty,4].
\]

\step{(b) Substitute the requested inputs into the appropriate formulas.}
All three inputs $0,1,2$ lie in both composition domains. Their
values, including the intermediate arithmetic, are
\[
 \begin{array}{c|c|c}
 x&(f\circ g)(x)=\sqrt{4-x^2}&(g\circ f)(x)=4-x\\ \hline
 0&\sqrt{4-0}=2&4-0=4\\
 1&\sqrt{4-1}=\sqrt3&4-1=3\\
 2&\sqrt{4-4}=0&4-2=2
 \end{array}
\]
Thus the six values in the order requested are
$2,4,\sqrt3,3,0,2$.

\step{(c) Check equality as functions, not only equality of formulas.}
The functions have different domains. Moreover, even on their
common domain they differ: at $0$ the first has value $2$ and
the second value $4$. Either observation disproves equality.

\step{(d) Test whether the inner value is admissible at $3$.}
Since $g(3)=9>4$, the expression $f(g(3))=f(9)$ is not defined
as a real value. In contrast, $3\le4$, so $f(3)=\sqrt1=1$
is defined and $g(f(3))=g(1)=1$ is defined.
\answer{The sum and product domains are $(-\infty,4]$; the two
composition domains are $[-2,2]$ and $(-\infty,4]$.
The compositions are unequal. At $3$, only $g\circ f$ is meaningful.}

\exercise{17.2}{Compositions can remove or preserve a jump}
\goal Let $f(x)=0$ for $x<0$, $f(x)=4$ for $x\ge0$, and
$g(x)=x^2$ for every real $x$. Determine the four combinations
in part (a), and decide which of the six functions in part (b)
are continuous.

\step{(a) Work separately on the two sides of zero.}
Both original domains are $\R$. Sums and products are therefore
defined on $\R$, with formulas
\[
 (f+g)(x)=\begin{cases}x^2,&x<0,\\x^2+4,&x\ge0,\end{cases}
 \qquad
 (fg)(x)=\begin{cases}0,&x<0,\\4x^2,&x\ge0.\end{cases}
\]
For every real $x$, $g(x)=x^2\ge0$. Thus $f(g(x))=4$,
including at $x=0$, because $f(0)=4$. Hence
\[
 (f\circ g)(x)=4\qquad(x\in\R).
\]
For the opposite order, square the two possible values of $f$:
\[
 (g\circ f)(x)=\begin{cases}0^2=0,&x<0,\\4^2=16,&x\ge0.\end{cases}
\]
The composition domains are $\R$ as well, since all inner
values belong to the outer domains.

\step{(b) First check points other than zero.}
At any $x_0\ne0$, a radius $\abs{x_0}/2$ keeps $x$ on the
same side of zero as $x_0$. Each displayed function is then
given near $x_0$ by one constant or polynomial formula.
Such formulas are continuous by the identity, constant, sum,
and product rules. Thus zero is the only possible discontinuity
of any function in this exercise.

\step{(b) Identify the three jumps with one sequence.}
Take $x_n=-1/n$. Then $x_n\to0$, but
\[
 f(x_n)=0\to0\ne4=f(0),
\]
\[
 (f+g)(x_n)=1/n^2\to0\ne4=(f+g)(0),
\]
and
\[
 (g\circ f)(x_n)=0\to0\ne16=(g\circ f)(0).
\]
Definition~17.1 therefore makes all three functions discontinuous
at zero.

\step{(b) Check the functions that remain continuous at zero.}
The function $g(x)=x^2$ is a product of two identity functions,
so it is continuous everywhere. The function $f\circ g$ is
constant, so it is continuous everywhere too.

For $fg$, its value at zero is $0$, and for every real $x$,
$\abs{(fg)(x)}\le4x^2$. Given $\varepsilon>0$, choose
$\delta=\sqrt{\varepsilon}/2>0$. If $\abs{x}<\delta$, then
\[
 \abs{(fg)(x)-(fg)(0)}\le4x^2<4\delta^2=\varepsilon.
\]
Hence $fg$ is continuous at zero and therefore on all of $\R$.
\answer{The globally continuous functions are $g$, $fg$, and
$f\circ g$. The functions $f$, $f+g$, and $g\circ f$ are
continuous away from zero but discontinuous at zero.}

\exercise{17.3}{Continuity from algebra and composition}
\goal Prove continuity of all eight requested functions using
the given familiar continuous functions and the section's rules.
\idea Decompose each expression into known continuous maps,
then check that every outer function receives an allowed input.

\step{(a) The function $\log(1+\cos^4x)$.}
The given continuity of cosine and repeated products show
that $\cos^4x$ is continuous on $\R$. Addition of the constant
$1$ preserves continuity. Furthermore,
$1+\cos^4x\ge1>0$ for every real $x$.
The logarithm is continuous on $(0,\infty)$, which contains
all these inner values. The composition theorem therefore proves
continuity on the natural domain $\R$.

\step{(b) The function $[\sin^2x+\cos^6x]^\pi$.}
Products and a sum show that $u(x)=\sin^2x+\cos^6x$ is
continuous on $\R$. Both summands are nonnegative.
They cannot both vanish: if $\sin x=0$, the identity
$\sin^2x+\cos^2x=1$ gives $\cos^2x=1$, so $\cos^6x=1$.
If $\sin x\ne0$, the first summand is positive instead.
Thus $u(x)>0$ for every $x$.
The given power function $t\mapsto t^\pi$ is continuous for
$t>0$, so its composition with $u$ is continuous on $\R$.

\step{(c) The function $2^{x^2}$.}
The map $x\mapsto x^2=I(x)I(x)$ is continuous on $\R$.
The given map $t\mapsto2^t$ is continuous for every real $t$.
Its composition with $x^2$ is therefore continuous on $\R$.

\step{(d) The function $8^x$.}
Use $8^x=2^{3x}$. The inner map $x\mapsto3x$ is continuous
as a scalar multiple of the identity, and the outer map
$t\mapsto2^t$ is continuous on $\R$. Composition proves the
claim on $\R$ without needing a new rule for the base $8$.

\step{(e) The tangent function on its stated domain.}
Write $\tan x=\sin x/\cos x$. Sine and cosine are continuous
on $\R$, and cosine is nonzero exactly away from
$x=\pi/2+k\pi$, $k\in\Z$. The quotient theorem proves
continuity at every point of that natural domain.

\step{(f) The function $x\sin(1/x)$ for $x\ne0$.}
The constant function $1$ and identity are continuous; the
quotient $r(x)=1/x$ is therefore continuous on
$D=\R\setminus\{0\}$. Since sine is continuous on all of
$\R$, $\sin(r(x))$ is continuous on $D$ by composition.
Multiplying by the continuous identity gives continuity of
$x\sin(1/x)$ on $D$.

\step{(g) The function $x^2\sin(1/x)$ for $x\ne0$.}
The sine factor is continuous on $D$ by part (f)'s construction.
The factor $x^2$ is continuous on $\R$, hence on $D$.
Their product is continuous at every point of $D$.

\step{(h) The function $x^{-1}\sin(1/x^2)$ for $x\ne0$.}
On $D$, $x^2\ne0$. The quotient $1/x^2$ is continuous there,
and its composition with sine is continuous there.
The other factor $1/x$ is continuous on $D$ too. Their product
is the requested continuous function.
\answer{Every function is continuous on its natural stated domain:
$\R$ in (a)--(d), the nonzero-cosine domain in (e), and
$\R\setminus\{0\}$ in (f)--(h). No value at zero is asserted
in the last three parts.}

\exercise{17.4}{Continuity of the square root including its endpoint}
\goal Prove that $f(x)=\sqrt x$ is continuous on $[0,\infty)$.
\idea An estimate valid even at zero avoids dividing by a
possibly small sum of square roots.

\step{Prove the estimate for any two domain points.}
If $x\ge y\ge0$, then $\sqrt x-\sqrt y\ge0$ and
$\sqrt x+\sqrt y\ge\sqrt x-\sqrt y$. Multiplication by
$\sqrt x-\sqrt y$ preserves this order, giving
\[
 x-y=(\sqrt x-\sqrt y)(\sqrt x+\sqrt y)
 \ge(\sqrt x-\sqrt y)^2.
\]
Taking nonnegative square roots yields
$\abs{\sqrt x-\sqrt y}\le\sqrt{\abs{x-y}}$.
If $y>x$, interchange the two variables; the same absolute-value
inequality results. Thus it holds for every $x,y\ge0$.

\step{Choose delta at an arbitrary point of the domain.}
Fix $x_0\ge0$ and $\varepsilon>0$, and put
$\delta=\varepsilon^2>0$. For a domain point $x\ge0$ with
$\abs{x-x_0}<\delta$, the estimate gives
\[
 \abs{\sqrt x-\sqrt{x_0}}
 \le\sqrt{\abs{x-x_0}}<\sqrt\delta=\varepsilon.
\]
This is Theorem~17.2. It includes $x_0=0$, because the condition
tests only $x\ge0$ and never requires a negative input.
\answer{The square-root function is continuous at every point of
its domain, including zero.}

\exercise{17.5}{Powers and polynomials are continuous}
\goal Prove part (a) for every positive integer power, then
deduce part (b) for every polynomial.

\step{(a) Build powers by induction using the product rule.}
For $m=1$, $x\mapsto x^m=x$ is the identity, already shown
continuous. Suppose $x\mapsto x^m$ is continuous on $\R$.
Then $x^{m+1}=x^m x$ is the product of that continuous function
and the identity. Theorem~17.4(ii) makes it continuous on $\R$.
Mathematical induction proves continuity for all $m\in\N$.

\step{(b) Assemble a polynomial from finitely many continuous terms.}
Let
\[
 p(x)=a_0+a_1x+\cdots+a_nx^n,
 \qquad a_0,\ldots,a_n\in\R.
\]
The term $a_0$ is constant and continuous. For $j\ge1$,
part (a) makes $x^j$ continuous, and Theorem~17.3 makes
$a_jx^j$ continuous, including when $a_j=0$.
Theorem~17.4(i) preserves continuity under addition of two
terms. Applying it finitely many times proves continuity of
the displayed sum at every real point. A constant polynomial
is already covered even if there are no positive-power terms.
\answer{Every positive integer power and every real polynomial
is continuous on $\R$.}

\exercise{17.6}{Continuity of rational functions on their domains}
\goal Let $p,q$ be polynomial functions and define
\[
 D=\{x\in\R:q(x)\ne0\},\qquad f(x)=p(x)/q(x)\quad(x\in D).
\]
Prove $f$ is continuous.

\step{Check the hypotheses at an arbitrary allowed point.}
Fix $x_0\in D$. Exercise~17.5 makes both $p$ and $q$
continuous at $x_0$, and membership in $D$ gives $q(x_0)\ne0$.
The quotient theorem, Theorem~17.4(iii), therefore applies and
gives continuity of $p/q$ at $x_0$.

For a direct reading of that theorem, take any $x_n\in D$
with $x_n\to x_0$. Continuity gives $p(x_n)\to p(x_0)$
and $q(x_n)\to q(x_0)\ne0$. The quotient limit law gives
\[
 f(x_n)=\frac{p(x_n)}{q(x_n)}\longrightarrow
 \frac{p(x_0)}{q(x_0)}=f(x_0).
\]
Every denominator in this display is nonzero by the domain
conditions. Since $x_0\in D$ was arbitrary, the function is
continuous on $D$.
\answer{A rational function is continuous wherever its denominator
does not vanish. Excluded zeros of the denominator are not
points at which continuity of this domain function is tested.}

\exercise{17.7}{Linear maps, absolute value, and scalar multiples}
\goal Prove (a) continuity of $x\mapsto kx$, (b) continuity of
absolute value, and (c) Theorem~17.3 by composition.

\step{(a) A linear map is a polynomial.}
For a fixed real $k$, the function $G(t)=kt$ is a polynomial
with constant coefficient zero and linear coefficient $k$.
Exercise~17.5 makes it continuous on $\R$. This includes
$k=0$, in which case it is constant.

\step{(b) Bound the error in absolute value by the input error.}
The triangle inequality gives
$\abs{x}\le\abs{x-x_0}+\abs{x_0}$, so
$\abs{x}-\abs{x_0}\le\abs{x-x_0}$.
Interchanging $x$ and $x_0$ gives
$\abs{x_0}-\abs{x}\le\abs{x-x_0}$.
Together these two inequalities are
\[
 \abs{\abs{x}-\abs{x_0}}\le\abs{x-x_0}.
\]
Given $\varepsilon>0$, choose $\delta=\varepsilon$.
Then $\abs{x-x_0}<\delta$ forces
$\abs{\abs{x}-\abs{x_0}}<\varepsilon$.
Thus $H(t)=\abs{t}$ is continuous at every real point.

\step{(c) Compose with the given function.}
Suppose $f$ is continuous at a point $x_0$ of its domain.
The outer maps $G$ and $H$ are continuous at the real number
$f(x_0)$ by parts (a) and (b). Theorem~17.5 gives continuity
at $x_0$ of
\[
 (G\circ f)(x)=kf(x),\qquad
 (H\circ f)(x)=\abs{f(x)}.
\]
The outer maps have domain $\R$, so both compositions retain
the whole domain of $f$. These are precisely the two functions
in Theorem~17.3.
\answer{Linear maps and absolute value are continuous on $\R$;
composition with them gives a second proof of Theorem~17.3.}

\exercise{17.8}{Formulas and continuity for a pointwise minimum}
\goal Verify both pointwise identities and use one to prove
the continuity assertion in part (c).

\step{(a) Verify the absolute-value formula by cases.}
For arbitrary real numbers $a,b$, suppose first that $a\le b$.
Then $\abs{a-b}=b-a$, and
\[
 \frac{a+b-\abs{a-b}}2
 =\frac{a+b-(b-a)}2=a=\min\{a,b\}.
\]
If $a>b$, then $\abs{a-b}=a-b$, and the same expression is
$(a+b-(a-b))/2=b=\min\{a,b\}$.
Substituting $a=f(x)$ and $b=g(x)$ wherever both are defined gives
\[
 \min(f,g)=\tfrac12(f+g)-\tfrac12\abs{f-g}.
\]

\step{(b) Verify the negated-maximum formula.}
If $a\le b$, negation reverses order, so $-a\ge-b$.
Thus $\max\{-a,-b\}=-a$ and its negative is $a$.
If $a>b$, one has $-a<-b$, so the negative of the maximum
is $b$. In both cases,
$\min\{a,b\}=-\max\{-a,-b\}$.
Pointwise substitution therefore gives
$\min(f,g)=-\max(-f,-g)$.

\step{(c) Apply the continuity rules to part (a).}
Suppose $f$ and $g$ are continuous at $x_0$ and both are
defined there. Their sum and difference are continuous at
$x_0$ by the sum and scalar rules. The absolute-value rule
makes $\abs{f-g}$ continuous there. Multiplying the two terms
of part (a) by $1/2$ and $-1/2$, then adding, preserves
continuity. The resulting function is $\min(f,g)$.
Its domain is the common domain of $f$ and $g$, so the point
$x_0$ is an allowed point throughout the argument.
\answer{Both identities hold, and the pointwise minimum of two
functions continuous at $x_0$ is continuous there.}

\exercise{17.9}{Four explicit epsilon--delta proofs}
\goal Verify the epsilon--delta condition directly in each
part, with a stated positive choice of $\delta$.

\step{(a) The square function at $x_0=2$.}
Fix $\varepsilon>0$. Factor the error:
\[
 \abs{x^2-4}=\abs{x-2}\abs{x+2}.
\]
To bound the second factor, first require $\abs{x-2}<1$.
Then
$\abs{x+2}\le\abs{x-2}+4<5$.
Choose
\[
 \delta=\min\{1,\varepsilon/5\}>0.
\]
If $\abs{x-2}<\delta$, both preliminary restrictions hold,
and $\abs{x^2-4}\le5\abs{x-2}<5\delta\le\varepsilon$.
This proves continuity at $2$.

\step{(b) The square-root function at zero.}
The domain restriction is $x\ge0$. Given $\varepsilon>0$,
put $\delta=\varepsilon^2>0$. If $x\ge0$ and
$\abs{x-0}<\delta$, then $0\le x<\varepsilon^2$ and
\[
 \abs{\sqrt x-\sqrt0}=\sqrt x<\varepsilon.
\]
The proof tests only domain points, as required at this endpoint.

\step{(c) The function $x\sin(1/x)$ extended by zero.}
Let $f(0)=0$ and $f(x)=x\sin(1/x)$ for $x\ne0$.
For nonzero $x$, the bound $\abs{\sin(1/x)}\le1$ gives
$\abs{f(x)}\le\abs{x}$. It also holds at $x=0$, where
both sides are zero. Given $\varepsilon>0$, choose
$\delta=\varepsilon$. If $\abs{x}<\delta$, then
\[
 \abs{f(x)-f(0)}=\abs{f(x)}\le\abs{x}<\varepsilon.
\]
Hence $f$ is continuous at zero despite the oscillating sine factor.

\step{(d) The cube function at an arbitrary $x_0\in\R$.}
Fix $\varepsilon>0$ and set $M=\abs{x_0}+1>0$.
If $\abs{x-x_0}<1$, then
$\abs{x}\le\abs{x-x_0}+\abs{x_0}<M$ and
$\abs{x_0}<M$. Thus
\[
 \abs{x^2+xx_0+x_0^2}
 \le\abs{x}^2+\abs{x}\abs{x_0}+\abs{x_0}^2
 \le3M^2.
\]
Using the hinted factorization, choose
\[
 \delta=\min\left\{1,\frac{\varepsilon}{3M^2}\right\}>0.
\]
For $\abs{x-x_0}<\delta$, the local bound applies, and
\[
 \abs{x^3-x_0^3}
 =\abs{x-x_0}\abs{x^2+xx_0+x_0^2}
 \le3M^2\abs{x-x_0}<3M^2\delta\le\varepsilon.
\]
\answer{The stated delta choices verify continuity in all four
parts directly, with all local bounds imposed before use.}

\exercise{17.10}{Three explicit sequences witnessing discontinuity}
\goal Show discontinuity at zero in parts (a)--(c) by finding
domain sequences approaching zero whose function values do not
approach the assigned value zero.

\step{(a) The one-sided step function.}
Here $f(x)=1$ for $x>0$ and $f(x)=0$ for $x\le0$.
Take $x_n=1/n$. Then $x_n>0$ and $x_n\to0$, but
$f(x_n)=1$ for all $n$. The image sequence tends to $1$,
whereas $f(0)=0$. This violates Definition~17.1.

\step{(b) The oscillating sine function.}
Here $g(x)=\sin(1/x)$ for $x\ne0$ and $g(0)=0$.
Put
\[
 x_n=\frac1{2\pi n+\pi/2}.
\]
The denominator is positive and tends to $+\infty$, so
$x_n>0$ and $x_n\to0$. At these inputs,
\[
 g(x_n)=\sin(2\pi n+\pi/2)=1.
\]
Thus $g(x_n)$ tends to $1$, not to $g(0)=0$.

\step{(c) The signum function.}
For $\operatorname{sgn}(x)=1$ at positive inputs, $-1$ at
negative inputs, and $0$ at zero, again take $x_n=1/n$.
Then $x_n\to0$ while $\operatorname{sgn}(x_n)=1$ for all
$n$, contrary to the image limit required by continuity.
\answer{Each function is discontinuous at zero. In every case
the displayed sequence stays in the domain and approaches the
point, while its image remains at distance $1$ from the value there.}

\exercise{17.11}{Testing continuity using only monotone sequences}
\goal Prove that $f$ is continuous at a domain point $x_0$
if and only if every monotone domain sequence $x_n\to x_0$
satisfies $f(x_n)\to f(x_0)$.
\idea For the reverse implication, a failure of continuity
produces a bad sequence; Theorem~11.4 supplies a monotone
subsequence of that sequence without repairing its bad images.

\step{Forward implication: restrict the sequences already tested.}
If $f$ is continuous at $x_0$, Definition~17.1 gives the image
limit for every domain sequence approaching $x_0$. In particular
it gives it for every monotone such sequence.

\step{Reverse implication: express failure of continuity precisely.}
Assume the stated condition for monotone sequences, and suppose
for contradiction that $f$ is not continuous at $x_0$.
The negation of Theorem~17.2 supplies one fixed
$\varepsilon_0>0$ such that, for every $\delta>0$, there is
$u$ in the domain with
\[
 \abs{u-x_0}<\delta,\qquad
 \abs{f(u)-f(x_0)}\ge\varepsilon_0.
\]
For each $n\in\N$, use $\delta=1/n$ and select such a point
$u_n$. Then $\abs{u_n-x_0}<1/n$ proves $u_n\to x_0$,
but every image has error at least $\varepsilon_0$.

\step{Select monotone inputs while retaining their bad images.}
Theorem~11.4 gives a monotone subsequence $(u_{n_k})$.
It still lies in the domain, and Theorem~11.3 gives
$u_{n_k}\to x_0$. The assumed monotone-sequence criterion
would therefore give $f(u_{n_k})\to f(x_0)$.
That is impossible, since
\[
 \abs{f(u_{n_k})-f(x_0)}\ge\varepsilon_0
 \qquad\text{for every }k.
\]
For example, no image can satisfy the eventual error bound
$\varepsilon_0/2$. The contradiction proves continuity.
\answer{Monotone test sequences suffice. The argument selects
a monotone subsequence; it does not assume the original sequence
of bad inputs is monotone.}

\exercise{17.12}{A continuous function is determined on a dense set}
\goal Prove (a) that zero values at all rational points of
$(a,b)$ force zero values everywhere, and (b) that equality of
two continuous functions at those rational points forces equality
throughout the interval.

\step{(a) Approximate an arbitrary interval point by rational points inside it.}
Fix $x\in(a,b)$. For each $n\in\N$, put
\[
 \rho_n=\min\{(x-a)/2,(b-x)/2,1/n\}>0.
\]
The open interval $(x-\rho_n,x+\rho_n)$ lies in $(a,b)$
and has positive length. Rational density supplies
$r_n\in\Q$ in that interval. In particular,
$\abs{r_n-x}<\rho_n\le1/n$, so $r_n\to x$.
Since $f$ is continuous at $x$ and $f(r_n)=0$ for every $n$,
\[
 f(x)=\lim_{n\to\infty}f(r_n)=\lim_{n\to\infty}0=0.
\]
The point $x$ was arbitrary, so $f$ vanishes everywhere on
$(a,b)$.

\step{(b) Reduce equality to the zero-valued case.}
The difference $h=f-g$ is continuous on $(a,b)$ by the sum
and scalar rules. At every rational point $r$ in that interval,
the hypothesis gives $h(r)=f(r)-g(r)=0$.
Part (a), applied to $h$, yields $h(x)=0$ for every $x\in(a,b)$.
Rearranging $f(x)-g(x)=0$ gives $f(x)=g(x)$ there.
\answer{Both assertions follow by taking limits along rational
points that stay inside the common domain.}

\exercise{17.13}{Rational and irrational inputs with different values}
\goal Prove (a) that the rational-indicator function is
discontinuous everywhere, and (b) that the function equal to
$x$ at rational inputs and zero at irrational inputs is continuous
exactly at zero.

\step{Preparation: obtain both kinds of inputs near any real point.}
Rationals are dense in $\R$. Irrationals are dense as well:
given $u<v$, rational density gives a rational $r$ in
$(u-\sqrt2,v-\sqrt2)$. Then $r+\sqrt2\in(u,v)$ is
irrational, since rationality of that sum would make $\sqrt2$
the difference of two rationals.

For each real $a$ and each $n$, choose a rational $r_n$ and
an irrational $t_n$ in $(a-1/n,a+1/n)$. Both sequences tend
to $a$, since their distances from $a$ are less than $1/n$.

\step{(a) The indicator takes incompatible limits along these inputs.}
Let $f(x)=1$ for rational $x$ and $f(x)=0$ for irrational $x$.
At the two sequences just chosen,
$f(r_n)=1$ and $f(t_n)=0$ for every $n$.
If $f$ were continuous at $a$, both image sequences would have
to tend to $f(a)$. Their limits are $1$ and $0$, which are
different. Thus $f$ cannot be continuous at $a$.
As $a$ was arbitrary, it is discontinuous everywhere.

\step{(b) At zero, the magnitude of the function is small with the input.}
Let $h(x)=x$ for rational $x$ and $h(x)=0$ for irrational $x$.
Then $h(0)=0$, and in both cases $\abs{h(x)}\le\abs{x}$.
Given $\varepsilon>0$, choose $\delta=\varepsilon$.
For $\abs{x}<\delta$,
\[
 \abs{h(x)-h(0)}\le\abs{x}<\varepsilon.
\]
Hence $h$ is continuous at zero.

\step{(b) At a nonzero point, the two dense input sets still disagree.}
Fix $a\ne0$ and choose the rational and irrational sequences
from the preparation step. Then
\[
 h(r_n)=r_n\longrightarrow a,\qquad
 h(t_n)=0\longrightarrow0.
\]
Continuity at $a$ would require both limits to equal $h(a)$.
Because $a\ne0$, the two limits differ, giving a contradiction.
\answer{The function in (a) is nowhere continuous. The function
in (b) is continuous at zero and at no other point.}

\exercise{17.14}{A denominator-sensitive function on the rationals}
\goal For nonzero rational $x=p/q$ in lowest terms with $q>0$,
let $f(x)=1/q$; also let $f(0)=1$ and let $f$ vanish on
irrationals. Prove continuity at every irrational point and
discontinuity at every rational point.
\idea Near an irrational point, only finitely many fractions
in a bounded neighborhood have small denominators. Avoid those
finitely many fractions; all remaining nonzero values of $f$
are small.

\step{Unify the notation at the rational point zero.}
Write $0=0/1$ in lowest terms. The same denominator rule
then gives $1/q=1=f(0)$. Thus all rational inputs, including
zero, can be discussed with a positive reduced denominator.

\step{Continuity at an irrational point: isolate the finitely many exceptions.}
Fix $a\in\R\setminus\Q$ and $\varepsilon>0$.
Choose an integer $Q>\max\{1,1/\varepsilon\}$, so
$1/Q<\varepsilon$. Consider
\[
 F=\left\{\frac pq\in[a-1,a+1]:
 p\in\Z,\ 1\le q\le Q,\ p/q\text{ is in lowest terms}\right\}.
\]
This set is finite. Indeed, for a fraction in this set,
\[
 \abs{p}=q\abs{p/q}\le q(\abs a+1)\le Q(\abs a+1).
\]
There are only finitely many possible integers $p$ in this
fixed bounded range, and only $Q$ possible denominators.
Their possible pairs, and hence the resulting fractions, are finite.
The argument concerns a bounded neighborhood; fractions with
small denominators are not a finite set on the whole real line.

No member of $F$ equals $a$, because $a$ is irrational.
If $F\ne\varnothing$, let
$d=\min\{\abs{r-a}:r\in F\}>0$ and set
$\delta=\min\{1,d/2\}$. If $F=\varnothing$, set
$\delta=1$ instead. In either case $\delta>0$.

\step{Check every nearby input, including rational ones.}
Suppose $\abs{x-a}<\delta$. Then $x\in[a-1,a+1]$,
and $x\notin F$ by the choice of $\delta$.
If $x$ is irrational, $f(x)=0=f(a)$.
If $x=p/q$ is rational in lowest terms, its denominator cannot
be at most $Q$, since that would put $x$ in $F$.
Thus $q>Q$ and
\[
 \abs{f(x)-f(a)}=\frac1q<\frac1Q<\varepsilon.
\]
The same strict error bound is immediate in the irrational case.
This proves continuity at $a$ by Theorem~17.2.

\step{Discontinuity at a rational point: approach through irrationals.}
Fix a rational $a=p/q$ in lowest terms, allowing $0=0/1$.
Then $f(a)=1/q>0$. By irrational density, choose irrational
$t_n\in(a-1/n,a+1/n)$ for each $n$.
These points tend to $a$, while $f(t_n)=0$ for every $n$.
The image limit is zero and differs from the positive value
$f(a)$, so $f$ is discontinuous at $a$.
\answer{The continuity set is exactly $\R\setminus\Q$;
the discontinuity set is exactly $\Q$.}

\exercise{17.15}{A continuity test that omits the point itself}
\goal Prove that $f$ is continuous at a domain point $x_0$
if and only if $f(x_n)\to f(x_0)$ for every sequence in
$\operatorname{dom}(f)\setminus\{x_0\}$ tending to $x_0$.

\step{Forward implication: the omitted-point sequences are a subclass.}
If $f$ is continuous at $x_0$, Definition~17.1 gives the image
limit for every domain sequence approaching $x_0$.
It therefore gives it for each sequence that avoids $x_0$.

\step{Reverse implication: split an arbitrary sequence by its equal terms.}
Assume the omitted-point condition. Let $(y_n)$ be any domain
sequence with $y_n\to x_0$. At indices where $y_n=x_0$,
the image error is exactly zero:
$\abs{f(y_n)-f(x_0)}=0$.
If only finitely many indices have $y_n\ne x_0$, the sequence
is eventually equal to $x_0$, so its images are eventually
equal to $f(x_0)$ and converge to that value.

If infinitely many indices have $y_n\ne x_0$, list all of
them in increasing order as $n_1<n_2<\cdots$.
The subsequence $x_k=y_{n_k}$ avoids $x_0$ and still tends
to $x_0$. The assumption gives $f(x_k)\to f(x_0)$.
Given $\varepsilon>0$, choose an integer $K\ge1$ so that
$k>K$ implies $\abs{f(x_k)-f(x_0)}<\varepsilon$.
For any original index $n>n_K$, either $y_n=x_0$, giving
zero error, or $n=n_k$ for some $k>K$, giving the same
strict error bound by the selected subsequence.
Thus $f(y_n)\to f(x_0)$ for the original full sequence.
As the sequence was arbitrary, Definition~17.1 proves continuity.

\step{Check what happens at an isolated domain point.}
If some radius $r>0$ satisfies
\[
 \operatorname{dom}(f)\cap(x_0-r,x_0+r)=\{x_0\},
\]
there is no sequence of domain points avoiding $x_0$ that
converges to $x_0$: convergence would eventually place its
points in this ball, where no such points exist.
The omitted-point condition is then vacuously true.
Every domain sequence converging to $x_0$ is eventually
$x_0$, so continuity is automatically true as well.
\answer{The punctured-sequence condition is equivalent to
continuity, including at isolated points of the domain.}

\exercise{17.16}{Jumps in the postage-stamp function}
\goal For fixed postage charges $A$ and $B$, with the additional
ounce charge $B>0$, prove that the stated function $P$ is
discontinuous at every positive integer.
\idea Approach an integer from the preceding interval. Its
charge differs from the value at the integer by exactly $B$.

\step{Write the domain and the two adjacent charges.}
The domain is $[0,\infty)$. The charge is $A$ on $[0,1)$,
and is $A+Bn$ on $[n,n+1)$ for every positive integer $n$.
Equivalently, the latter formula also covers $n=0$.
Fix any positive integer $m$. Its value is
$P(m)=A+Bm$, because $m\in[m,m+1)$.
On the interval $[m-1,m)$, the value is instead
$A+B(m-1)$.

\step{Select domain points in the preceding interval.}
For each $n\in\N$, set
\[
 x_n=m-\frac1{n+1}.
\]
Since $0<1/(n+1)<1$, we have $m-1<x_n<m$, so $x_n$
belongs to the domain and to the preceding charge interval.
Also $\abs{x_n-m}=1/(n+1)\to0$, hence $x_n\to m$.
But
\[
 P(x_n)=A+B(m-1)\longrightarrow A+B(m-1)\ne A+Bm=P(m).
\]
The inequality holds because their difference is the nonzero
charge $B$. Definition~17.1 therefore gives discontinuity at $m$.
Since $m$ was any positive integer, all such integers are
discontinuity points.

\step{Record the implicit nonzero-charge assumption.}
The proof works for any fixed $B\ne0$, although postage gives
the natural assumption $B>0$. If $B=0$, the formula reduces
to the constant function $P(x)=A$, which is continuous;
the requested jump assertion necessarily presumes a nonzero
additional-ounce charge. Charges here are fixed for the rate
schedule under consideration.
\answer{With $B\ne0$, each positive integer has a jump of
magnitude $\abs B$, so $P$ is discontinuous at each of them.}
